\documentclass[11pt]{beamer}
\usecolortheme{dolphin}
\usetheme{Madrid}
\usefonttheme{serif} 
\usepackage{amsmath, amssymb, graphicx, booktabs}
\usepackage{siunitx}
\usepackage{bm}
\usepackage{multirow}
\usepackage{caption}
\usepackage{mathptmx}
\captionsetup[figure]{font=scriptsize,labelfont=bf}

% ============================================================
% ADD LOGO TO FOOTLINE
% ============================================================

\setbeamertemplate{footline}{
	\leavevmode
	\hbox{%
		% Left box: footer text
		\begin{beamercolorbox}[wd=.68\paperwidth,ht=2.5ex,dp=1ex,left]{author in head/foot}
			\usebeamerfont{author in head/foot}\hspace{1em}
			Time Series Analysis Using Ordinal Patterns
		\end{beamercolorbox}%
		
		% Middle box: slide numbers
		\begin{beamercolorbox}[wd=.17\paperwidth,ht=2.5ex,dp=1ex,center]{date in head/foot}
			\usebeamerfont{date in head/foot}
			\insertframenumber{} / \inserttotalframenumber
		\end{beamercolorbox}%
		
		% Right box: logo
		\begin{beamercolorbox}[wd=.15\paperwidth,ht=2.5ex,dp=1ex,right]{date in head/foot}
			\raisebox{-0.2ex}{
			%	
\includegraphics[height=2ex]{"C:/Users/UserA1/Documents/GitHub/Ordinal_Patterns_R/Logos/vuw_logo.png"}
			}
			\hspace{1em}
		\end{beamercolorbox}%
	}
	\vskip0pt
}

% ============================================================

\title[Ordinal Pattern Analysis of Time Series]{Feature-Based Clustering of Time Series \\ Using Ordinal Pattern Analysis}
\author[M. H. M. R. S. Dilhani et al.]{
	M.\ H.\ M.\ R.\ S.\ Dilhani\inst{1} \\
	Alejandro C.\ Frery\inst{1} \\
	Andrea A.\ Rey\inst{2}
}
\institute[Victoria University of Wellington]{
	\inst{1}School of Mathematics and Statistics, Victoria University of Wellington, New Zealand\\
	\inst{2}LIDEC, Universidad Nacional de Hurlingham (UNAHUR), Argentina
}

\date{08 December 2025}
\titlegraphic{%
	
\includegraphics[width=4cm]{Logo Offshore Standard Landscape Reversed RGB.png}\hfill
	
\includegraphics[width=4cm]{NZSA-logo-words-e1505785105112-1.png}%
}

\begin{document}
	
	\begin{frame}
		\titlepage
	\end{frame}
	
	\begin{frame}{Outline}
		\tableofcontents
	\end{frame}
	
	% ============================================================
	% YOUR REMAINING SLIDES — UNCHANGED
	% ============================================================
	
	\begin{frame}{Motivation and Background}
		\small
		\begin{itemize}
			\item Time series clustering is challenging since every series analyzed contains a stochastic component, contributing random variability and increased complexity.
			\item \alert{Ordinal patterns} capture the order relations among consecutive data points.
			\item \textbf{Bandt \& Pompe (2002)}~\cite{PhysRevLett.88.174102}: introduced \textbf{Permutation Entropy}.
			\item \textbf{Limitation:} Entropy quantifies disorder but cannot distinguish random from structured dynamics.
			\item \textbf{Extension:} \alert{Statistical Complexity}.
			\item The \textbf{Entropy–Complexity Plane} maps time series by entropy $(H)$ and complexity $(C)$.
			\item \textbf{Goal:} Differentiate AR, MA, and ARMA models.
		\end{itemize}
	\end{frame}
	
%-----------------------------------------------------%
\begin{frame}{From Time Series to Ordinal Patterns: An Example}
	\small
	\begin{itemize}
		\item \textbf{Example:} Mean monthly humidity in Wellington.
		\item Original time series:
		\[
		77.3,\ 81.0,\ 82.4,\ 81.7,\ 83.6,\ 85.6,\ \dots
		\]
	\end{itemize}
	
	\begin{figure}[hbt]
		\centering
		\includegraphics[width=0.5\textwidth]{C:/Users/UserA1/Documents/GitHub/Ordinal_Patterns_R/Plots/Auckland & Dunedin Conference related plots/humidity graph} 
		\caption{Mean monthly humidity in Wellington}
		\label{fig:humidity}
	\end{figure}
\end{frame}

%-----------------------------------------------------%
\begin{frame}{Transforming into Ordinal Patterns}
	\small
	\begin{itemize}
		\item Choose embedding dimension $D = 3$.
		\item Slide a window of size $D$ across the series:
	\end{itemize}
	
	\[
	(77.3, 81.0, 82.4) \rightarrow (1,2,3) = \pi_1
	\]
	\[
	(81.0, 82.4, 81.7) \rightarrow (1,3,2) = \pi_2
	\]
	
	\begin{itemize}
		\item Continue for all overlapping windows ($n - D + 1$ total).
		\item Each window represents an \alert{ordinal pattern} describing local order relations.
	\end{itemize}
	
	\begin{center}
		\includegraphics[width=0.5\textwidth]{C:/Users/UserA1/Documents/GitHub/Ordinal_Patterns_R/Plots/Auckland & Dunedin Conference related plots/frequency histogram} 
	\end{center}
\end{frame}

%-----------------------------------------------------%
\begin{frame}
	
	\begin{block}{Definition: Permutation Entropy ($H$)}
		\[
		H(\mathbf{p}) = -\dfrac{1}{\log k} \sum_{i=1}^{k} p(\pi^i) \log p(\pi^i),
		\]
		where $k=D!$ and $p(\pi^i)$ is the proportion of pattern $\pi^i$.
	\end{block}
	
	\begin{block}{Definition: Jensen-Shannon Distance}
		Given a histogram $\mathbf{p}$ and the reference uniform distribution $\mathbf{u} = (1/k, \ldots, 1/k)$, the disequilibrium is calculated as:
		\[
		Q'(\mathbf{p},\mathbf{u}) = \sum_{\ell=1}^k \Big( p_\ell \log\frac{p_\ell}{u_\ell} + u_\ell \log\frac{u_\ell}{p_\ell} \Big).
		\]
	\end{block}		
	
\end{frame}



%-----------------------------------------------------%
\begin{frame}{Disequilibrium and Statistical Complexity}
	\begin{block}{Disequilibrium}
		The \textbf{normalized disequilibrium}:
		\[
		Q = \frac{Q'}{\max(Q')},
		\]
		provides a scale-independent comparison, where $\max(Q')$ is the highest possible value for the given $k$.
	\end{block}
	\begin{block}{Definition: Statistical Complexity $(C)$ }
		\[
		C = H Q
		\]
		where $H$ is normalized entropy and $Q$ is normalized disequilibrium.
	\end{block}
\end{frame}

%-----------------------------------------------------%
\begin{frame}
	Conceptual Framework
	\begin{columns}
		\column{0.55\textwidth}
		\begin{itemize}
			\item Transform time series $\{x_t\}$ into \textbf{ordinal patterns}.
			\item Compute:
			\begin{itemize}
				\item Permutation Entropy
				\item Statistical Complexity
				\item Confidence Interval for permutation entropy and complexity
			\end{itemize}
			\item Plot in the \textbf{Entropy–Complexity ($H \times C$)} plane.
			\item Use cluster formation to discriminate models.
		\end{itemize}
		\column{0.45\textwidth}
		\includegraphics[width=\linewidth]{C:/Users/UserA1/Documents/GitHub/Ordinal_Patterns_R/Plots/Auckland & Dunedin Conference related plots/New_model_group_plot_n1000}
		\captionof{figure}{ Entropy–Complexity plane for Sample size $n=1000$.}
	\end{columns}
\end{frame}

%===========================================================
\section{Methodology}
\begin{frame}{Experimental Design}
	\begin{itemize}
		\item Models studied:
		\begin{itemize}
			\item $\mathrm{AR}(1)$, $\mathrm{AR}(2)$, $\mathrm{MA}(1)$,  $\mathrm{MA}(2)$, $\mathrm{ARMA}(1,1)$, $\mathrm{ARMA}(2,2)$
		\end{itemize}
		\item Each model simulated with 100 independent replications.
		\item Parameter sets chosen to ensure stationarity and invertibility.
		\item Series lengths: $n = 1000, n=5000$.
	\end{itemize}
\end{frame}

%-----------------------------------------------------%
\begin{frame}{Model Parameter Structure}
	\centering
	\scriptsize
	\begin{tabular}{llrrrr}
		\toprule
		Type & Model & $\alpha_1$ & $\alpha_2$ & $\beta_1$ & $\beta_2$ \\
		\midrule
		$\mathrm{AR}(1)$ & $\mathrm{AR}(1)_{\textrm{M}_1}$ & $0.8$ &  &  & \\
		& $\mathrm{AR}(1)_{\textrm{M}_2}$ & $0.1$ &  &  & \\
		$\mathrm{MA}(1)$ & $\mathrm{MA}(1)_{\textrm{M}_1}$ &  &  & $0.8$ & \\
		& $\mathrm{MA}(1)_{\textrm{M}_2}$ &  &  & $0.1$ & \\
		$\mathrm{ARMA}(1,1)$ & $\mathrm{ARMA}(1,1)_{\textrm{M}_1}$ & $0.8$ &  & $0.8$ & \\
		& $\mathrm{ARMA}(1,1)_{\textrm{M}_2}$ & $0.1$ &  & $0.1$ & \\
		\bottomrule
	\end{tabular}
	\captionof{table}{Example of parameter sets ensuring diversity.}
\end{frame}

%-----------------------------------------------------%
\begin{frame}{Feature Extraction and Central Points}
	\begin{itemize}
			\item For each realization, compute:
			\begin{itemize}
				\item Entropy (\(H\)) — randomness
				\item Complexity (\(C\)) — structure
				\item Confidence Interval(CI) for $H$ and $C$
			\end{itemize}
			\item For each model, identify the \textbf{central point} $(H^*, C^*)$ using PCA projection.
			\item Central ensures the point is the most representative of the model’s behavior, as referred to in Chagas et al.~\cite{Chagas2022a}. 
		\end{itemize}
\end{frame}

%===========================================================
\section{Results: Case-Based Analysis}
\begin{frame}{Case Study Design}
	\begin{itemize}
		\item \textbf{Case 1–3:} Positive, negative, and mixed-sign coefficients.
		\item Two sample sizes: $n=1000, n=5000$.
		\item Each case analyzed using:
		\begin{itemize}
			\item The density of points in the $H \times C$
			\item Scatter plots with confidence intervals for $H$ and $C$ 
			\item Heat maps for model classification
		\end{itemize}
	\end{itemize}
\end{frame}

%-----------------------------------------------------% 	
\begin{frame}{Case 1 — Positive Coefficients}
		\begin{columns}[c] % 'c' centers content vertically
		\column{0.48\textwidth}
		\centering
		\includegraphics[width=\textwidth]{C:/Users/UserA1/Documents/GitHub/Ordinal_Patterns_R/Plots/Auckland & Dunedin Conference related plots/Scatter_Case1_n1000.pdf}
	%\captionof{figure}{Entropy–complexity scatter plot ($n=1000$) for positive model coefficients.}
		\column{0.48\textwidth}
	\centering
	\includegraphics[width=\textwidth]{C:/Users/UserA1/Documents/GitHub/Ordinal_Patterns_R/Plots/Auckland & Dunedin Conference related plots/Scatter_Case1_n5000.pdf}
	%\captionof{figure}{Entropy–complexity scatter plot ($n=5000$) for positive model coefficients.}
	\end{columns}
	\captionof{figure}{Entropy–complexity scatter plot ($n=1000$ and $n=5000$) for positive model coefficients.}
\end{frame}

\begin{frame}{Case 1 — Positive Coefficients cont\dots}
	\begin{columns}[c] % 'c' centers content vertically
		\column{0.48\textwidth}
		\centering
		\includegraphics[width=\textwidth]{C:/Users/UserA1/Documents/GitHub/Ordinal_Patterns_R/Plots/Auckland & Dunedin Conference related plots/CentralPoints_TimeSeries_Case1_n5000.pdf}
		%\captionof{figure}{Entropy–complexity scatter plot ($n=1000$) for positive model coefficients.}
		\column{0.48\textwidth}
		\centering
		\includegraphics[width=\textwidth]{C:/Users/UserA1/Documents/GitHub/Ordinal_Patterns_R/Plots/Auckland & Dunedin Conference related plots/HC_Heatmap_ModelTypes_Case1_n5000.pdf}
		%\captionof{figure}{Entropy–complexity scatter plot ($n=5000$) for positive model coefficients.}
	\end{columns}
	\captionof{figure}{Entropy–complexity central points and associated time series ($n=5000$) and heat map for positive model coefficients.}
\end{frame}

%\begin{frame}{Case 1 — Positive Coefficients cont\dots}
	%\begin{columns}[c] % 'c' centers content vertically
	%	\column{0.48\textwidth}
	%	\centering
	%	\includegraphics[width=\textwidth]{C:/Users/UserA1/Documents/GitHub/Ordinal_Patterns_R/Plots/Auckland & Dunedin Conference related plots/HC_Heatmap_ModelTypes_Case1_n1000.pdf}
		%\captionof{figure}{Entropy–complexity scatter plot ($n=1000$) for positive model coefficients.}
	%	\column{0.48\textwidth}
	%	\centering
	%	\includegraphics[width=\textwidth]{C:/Users/UserA1/Documents/GitHub/Ordinal_Patterns_R/Plots/Auckland & Dunedin Conference related plots/HC_Heatmap_ModelTypes_Case1_n5000.pdf}
		%\captionof{figure}{Entropy–complexity scatter plot ($n=5000$) for positive model coefficients.}
	%\end{columns}
	%\captionof{figure}{Entropy–complexity central points ($n=1000, n=5000$) for positive model coefficients.}
%\end{frame}

%-----------------------------------------------------%
\begin{frame}{Case 2 — Negative Coefficients}
\begin{columns}[c] % 'c' centers content vertically
	\column{0.48\textwidth}
	\centering
	\includegraphics[width=\textwidth]{C:/Users/UserA1/Documents/GitHub/Ordinal_Patterns_R/Plots/Auckland & Dunedin Conference related plots/Scatter_Case2_n1000.pdf}
	\column{0.48\textwidth}
	\centering
	\includegraphics[width=\textwidth]{C:/Users/UserA1/Documents/GitHub/Ordinal_Patterns_R/Plots/Auckland & Dunedin Conference related plots/Scatter_Case2_n5000.pdf}
\end{columns}
\captionof{figure}{Entropy–complexity scatter plot ($n=1000$ and $n=5000$) for negative model coefficients.}
\end{frame}

\begin{frame}{Case 2 — Negative Coefficients cont\dots}
	\begin{columns}[c] % 'c' centers content vertically
		\column{0.48\textwidth}
		\centering
		\includegraphics[width=\textwidth]{C:/Users/UserA1/Documents/GitHub/Ordinal_Patterns_R/Plots/Auckland & Dunedin Conference related plots/CentralPoints_TimeSeries_Case2_n5000.pdf}
		%\captionof{figure}{Entropy–complexity scatter plot ($n=1000$) for positive model coefficients.}
		\column{0.48\textwidth}
		\centering
		\includegraphics[width=\textwidth]{C:/Users/UserA1/Documents/GitHub/Ordinal_Patterns_R/Plots/Auckland & Dunedin Conference related plots/HC_Heatmap_ModelTypes_Case2_n5000.pdf}
		%\captionof{figure}{Entropy–complexity scatter plot ($n=5000$) for positive model coefficients.}
	\end{columns}
	\captionof{figure}{Entropy–complexity central points and associated time series ($n=5000$) and heat map for negative model coefficients.}
\end{frame}

%-----------------------------------------------------%
\begin{frame}{Case 3 — mixed-Sign Coefficients}
	\begin{columns}[c] % 'c' centers content vertically
		\column{0.48\textwidth}
		\centering
		\includegraphics[width=\textwidth]{C:/Users/UserA1/Documents/GitHub/Ordinal_Patterns_R/Plots/Auckland & Dunedin Conference related plots/Scatter_Case3_n1000.pdf}
		\column{0.48\textwidth}
		\centering
		\includegraphics[width=\textwidth]{C:/Users/UserA1/Documents/GitHub/Ordinal_Patterns_R/Plots/Auckland & Dunedin Conference related plots/Scatter_Case3_n5000.pdf}
	\end{columns}
	\captionof{figure}{Entropy–complexity scatter plot ($n=1000$ and $n=5000$) for mixed model coefficients.}
\end{frame}


\begin{frame}{Case 3 — Mixed-Sign Coefficients  cont\dots}
	\begin{columns}[c] % 'c' centers content vertically
		\column{0.48\textwidth}
		\centering
		\includegraphics[width=\textwidth]{C:/Users/UserA1/Documents/GitHub/Ordinal_Patterns_R/Plots/Auckland & Dunedin Conference related plots/CentralPoints_TimeSeries_Case3_n5000.pdf}
		%\captionof{figure}{Entropy–complexity scatter plot ($n=1000$) for positive model coefficients.}
		\column{0.48\textwidth}
		\centering
		\includegraphics[width=\textwidth]{C:/Users/UserA1/Documents/GitHub/Ordinal_Patterns_R/Plots/Auckland & Dunedin Conference related plots/HC_Heatmap_ModelTypes_Case3_n5000.pdf}
		%\captionof{figure}{Entropy–complexity scatter plot ($n=5000$) for positive model coefficients.}
	\end{columns}
	\captionof{figure}{Entropy–complexity central points and associated time series ($n=5000$) and heat map for mixed model coefficients.}
\end{frame}


%===========================================================
\section{Combined Model Visualization}
\begin{frame}{Combined Model Visualization}
	\centering
	\includegraphics[width=\textwidth]{C:/Users/UserA1/Documents/GitHub/Ordinal_Patterns_R/Plots/Auckland & Dunedin Conference related plots/combined_new_scatterplot.pdf}
	\captionof{figure}{All models combined in the entropy–complexity plane.}
	\centering
\end{frame}


%===========================================================
\section{Summary and Conclusion}
\begin{frame}{Key Findings}
	\begin{itemize}
		\item Ordinal pattern analysis successfully distinguishes AR, MA, and ARMA models.
		\item Distinct clustering appears across Permutation entropy.
		\item Central $(H^*, C^*)$ points provide consistent model representation.
		\item Larger samples yield tighter clusters and improved discrimination.
	\end{itemize}
\end{frame}

%-----------------------------------------------------%
\begin{frame}{Conclusion}
	\begin{block}{Summary}
		This study reveals that combining ordinal pattern analysis with entropy–complexity measures provides a robust framework for model discrimination in simulated time series.
	\end{block}
	\begin{itemize}
		\item Captures both linear and nonlinear relationship between the entropy and complexity.
		\item Overlap points reveal intrinsic similarities in model dynamics.
		\item Extensible to real-world applications such as finance, engineering, biomedical signals.
	\end{itemize}
\end{frame}

%-----------------------------------------------------%
\begin{frame}{Future Work}
	\begin{itemize}
		\item Apply unsupervised clustering (e.g., k-means, hierarchical) in the $H–C$ plane.
		\item Extend analysis by increasing sample size
		\item Extend distinct clustering using other types of entropies Rényi, Tsallis, and Fisher Information. 
		\item Extend analysis to empirical systems.
	\end{itemize}
\end{frame}

%-----------------------------------------------------%
\begin{frame}[plain]
	\centering
	\Huge \textbf{Thank You!} \\
	\vspace{0.5cm}
	\small Contact: \href{mailto:rasika.dilhani@vuw.ac.nz}{rasika.dilhani@vuw.ac.nz}
\end{frame}

%-----------------------------------------------------%
\begin{frame}[allowframebreaks]{References}
	\bibliographystyle{alpha}
	\bibliography{../../../../BearingFaultDiagnosis}
	
\end{frame}


	
\end{document}
